\subsection{Conjugate Gradient}

\begin{frame}{Conjugate Gradient: Explicación Gráfica}
    \begin{block}{Concepto principal}
        [Explica aquí la idea central de direcciones conjugadas y ortogonalidad respecto a la matriz hessiana].
    \end{block}
    \begin{itemize}
        \item {[Idea gráfica 1: Comportamiento en curvas de nivel cuadráticas frente al descenso estándar].}
        \item {[Idea gráfica 2: Ortogonalidad de las direcciones de búsqueda sucesivas].}
    \end{itemize}
\end{frame}

\begin{frame}{Conjugate Gradient: Pseudocódigo}
    \begin{algorithmic}[1]
        \State \textbf{Input:} {[Definir entradas: función \(f\), gradiente \(\nabla f\), punto inicial \(x^{(1)}\), tolerancias]}
        \State Inicializar dirección de descenso inicial \(d^{(1)}\)
        \While{[Condición de parada]}
            \State Calcular longitud de paso \(\alpha\)
            \State Actualizar punto: \(x \leftarrow x + \alpha d\)
            \State Calcular factor de conjugación \(\beta\)
            \State Actualizar nueva dirección: \(d \leftarrow -\nabla f + \beta d\)
        \EndWhile
        \State \textbf{Return} {[Punto óptimo]}
    \end{algorithmic}
\end{frame}

\begin{frame}{Conjugate Gradient: Ejemplo Paso a Paso}
    \begin{exampleblock}{Planteamiento del problema}
        [Define aquí una función de prueba cuadrática simple, ej: \(f(x_1, x_2) = x_1^2 + 2x_2^2\), y punto inicial].
    \end{exampleblock}
    \begin{itemize}
        \item \textbf{Paso 1:} {[Cálculo del gradiente inicial y primera dirección de búsqueda].}
        \item \textbf{Paso 2:} {[Cálculo del factor de actualización \(\beta\) y siguiente dirección].}
        \item \textbf{Paso 3:} {[Estado de convergencia y resultado].}
    \end{itemize}
\end{frame}

\begin{frame}[fragile]{Conjugate Gradient: Código}
\begin{lstlisting}[language=Python]
def conjugate_gradient(f, grad_f, x0, max_iter=50, tol=1e-5):
    # TODO: Inicializar variables y primera direccion d
    # while iter < max_iter:
    #     alpha = line_search(...)
    #     x = x + alpha * d
    #     beta = compute_beta(...)
    #     d = -grad + beta * d
    pass
\end{lstlisting}
\end{frame}

\begin{frame}{Conjugate Gradient: Análisis de Convergencia}
    \begin{alertblock}{Tasa de Convergencia}
        [Indicar tipo de convergencia en funciones cuadráticas vs funciones no lineales].
    \end{alertblock}
    \begin{itemize}
        \item {[Propiedades teóricas de convergencia en \(n\) dimensiones].}
        \item {[Ventajas y limitaciones frente a métodos de segundo orden puros].}
    \end{itemize}
\end{frame}